\begin{tikzpicture}[
  node distance=10mm and 2.5mm,
  every node/.style={font=\small},
  flow tag/.style={font=\footnotesize, fill=white, inner sep=1.2pt},
  input box/.style={flow box, minimum width=30mm, minimum height=10mm},
  skill box/.style={flow box, minimum width=37mm, minimum height=13mm},
  terminal box/.style={
    flow box,
    rounded corners=2mm,
    minimum width=38mm,
    minimum height=12mm
  }
]

% 第一层：异常恢复技能选择依据
\node[input box] (atype) {异常类型};
\node[input box, right=of atype] (task) {当前任务};
\node[input box, right=of task] (loc) {异常位置};

\coordinate (inputbus) at ($(task.south)+(0,-5mm)$);
\foreach \n in {atype,loc} {
  \draw[line width=0.6pt] (\n.south) -- (\n.south |- inputbus);
}
\draw[line width=0.6pt] (task.south) -- (task.south |- inputbus);
\draw[line width=0.6pt] (atype.south |- inputbus) --
  (loc.south |- inputbus);

\node[flow box, minimum width=48mm, below=10mm of task] (select)
  {恢复技能选择器};
\draw[flow arrow] (inputbus) -- (select.north);

% 第二层：三种恢复处理
\node[skill box, below left=16mm and 4mm of select] (vs)
  {视觉伺服恢复};
\node[skill box, below=16mm of select] (act)
  {动作块转换器恢复};
\node[terminal box, below right=16mm and 4mm of select] (manual)
  {人工接管};

\coordinate (skillbus) at ($(select.south)+(0,-6mm)$);
\draw[line width=0.6pt] (select.south) -- (skillbus);
\draw[line width=0.6pt] (vs.north |- skillbus) --
  (manual.north |- skillbus);
\draw[flow arrow] (vs.north |- skillbus) -- (vs.north);
\draw[flow arrow] (act.north |- skillbus) -- (act.north);
\draw[flow arrow] (manual.north |- skillbus) -- (manual.north);

\node[flow tag, align=center, anchor=south]
  at ($(vs.north |- skillbus)+(0,1mm)$)
  {进度停滞或已训练异常};
\node[flow tag, align=center, anchor=south]
  at ($(manual.north |- skillbus)+(0,1mm)$)
  {手术野未充分暴露\\或未知异常};

% 第三层：自主恢复汇聚
\node[flow box, minimum width=45mm, below=14mm of act] (execute)
  {执行恢复技能};
\draw[flow arrow] (vs.south) |- (execute.west);
\draw[flow arrow] (act.south) -- (execute.north);

% 恢复完成与异常事件判断
\node[flow decision, below=8mm of execute, minimum width=28mm] (reach)
  {恢复动作完成?\\[1pt]$\mathrm{Reach}(p_r)$?};
\draw[flow arrow] (execute) -- (reach);

\draw[flow arrow] (reach.west) -- ++(-7mm,0) |- (execute.west);
\node[flow tag, anchor=south]
  at ($(reach.west)+(-3.5mm,1mm)$) {否};

\node[flow decision, below=8mm of reach, minimum width=36mm] (abn)
  {异常事件仍成立?\\[1pt]$\mathrm{Abn}(p_{cv},p_{un},p_{prg})$?};
\draw[flow arrow] (reach) -- (abn);
\node[flow tag, anchor=east]
  at ($(reach.south)!0.50!(abn.north)+(-1mm,0)$) {是};

% 第四层：返回、重试或接管
\node[terminal box, below=10mm of abn] (return)
  {返回原正常任务节点};
\node[flow decision, right=4mm of abn,
  minimum width=24mm] (count)
  {恢复次数\\达到 $M$?};

\draw[flow arrow] (abn.south) -- (return.north);
\node[flow tag, anchor=south]
  at ($(abn.south)!0.50!(return.north)+(-2mm,0)$) {否};
\draw[flow arrow] (abn.east) -- (count.west);
\node[flow tag, anchor=south]
  at ($(abn.east)!0.50!(count.west)+(0,1mm)$) {是};

\coordinate (manualroute) at ($(manual.east)+(3mm,0)$);
\draw[flow arrow] (count.east) -- (manualroute |- count.east) --
  (manualroute) -- (manual.east);
\node[flow tag, anchor=south]
  at ($(count.east)+(5mm,1mm)$) {是};

\coordinate (looproute) at ($(vs.west)+(-3mm,0)$);
\coordinate (loopy) at ($(return.south)+(0,-5mm)$);
\draw[flow arrow] (count.south) -- (count.south |- loopy) --
  (looproute |- loopy) -- (looproute |- select.west) -- (select.west);
\node[flow tag, anchor=north]
  at ($(count.south)+(0,-1mm)$) {否，继续恢复};

\end{tikzpicture}
